% !TeX root = ../main.tex
\section{Vectorization Verification}\label{sec:vectorization}
\rubric{Vectorization verification}{15}{Clear evidence that SIMD instructions are used, with a correct
reading of vector width, FMA and masking.}

\begin{hypothesis}
  The \code{scalar} build should emit only scalar SSE (\code{ss}/\code{sd}) instructions and report
  zero vectorized loops. \code{autovec}, limited to the generic SSE2 target, should vectorize the
  streaming kernels at \qty{128}{bit} (\num{4} float32 lanes). \code{native} should vectorize them at
  \qty{256}{bit} (\code{ymm}, \num{8} float32 / \num{4} float64 lanes) and fuse the multiply--add in
  SAXPY and the stencil into \code{vfmadd}. We expect the compiler to \emph{refuse} to vectorize the
  dot-product reduction usefully, because summation reassociation is illegal without
  \code{-ffast-math}.
\end{hypothesis}

\Cref{tab:vec-summary} counts, per kernel function and variant, the instructions that use vector
registers or packed operations (\code{objdump}, \code{make report}). The few matches in the scalar
column are register-clearing \code{pxor}/\code{movaps}; GCC reports \emph{zero} vectorized loops
for \code{scalar}.

\begin{table}[tbp]
  \centering\small
  \caption{SIMD instructions per kernel function and build variant.}\label{tab:vec-summary}
  \IfDataExists{vectorization}{%
    \pgfplotstableset{acsvec/.style={
      col sep=comma,
      columns={variant,function,insns,simd_insns,widest_register},
      columns/variant/.style={column name=Build, string type},
      columns/function/.style={column name=Function, string type, string replace*={_}{\_}},
      columns/insns/.style={column name=Insns},
      columns/simd_insns/.style={column name=SIMD},
      columns/widest_register/.style={column name=Widest, string type},
      every head row/.style={before row=\toprule, after row=\midrule},
      every last row/.style={after row=\bottomrule}}}%
    \scriptsize\setlength{\tabcolsep}{2.5pt}%
    \pgfplotstabletypeset[acsvec, skip rows between index={12}{36}]{data/vectorization.csv}\hfill
    \pgfplotstabletypeset[acsvec, skip rows between index={0}{12}, skip rows between index={24}{36}]{data/vectorization.csv}\hfill
    \pgfplotstabletypeset[acsvec, skip rows between index={0}{24}]{data/vectorization.csv}%
  }
\end{table}

\subsection{Compiler report}
\code{make report} (\code{data/raw/vectorization/20261021-001032}) recompiles the unchanged
\code{src/kernels.c} with each variant's flags; \code{objdump} confirms the kernels match the timed
binaries instruction for instruction. \code{-fopt-info-vec} for SAXPY (\code{kernels.c:29}):
\begin{lstlisting}[style=excerpt]
# native  (-O3 -march=native):
src/kernels.c:29:26: optimized: loop vectorized using 32 byte vectors and unroll factor 8
src/kernels.c:29:26: optimized: epilogue loop vectorized using 16 byte vectors and unroll factor 4
# autovec (-O3, generic x86-64):
src/kernels.c:29:26: optimized: loop vectorized using 16 byte vectors and unroll factor 4
src/kernels.c:29:26: optimized: epilogue loop vectorized using 8 byte vectors and unroll factor 2
# scalar  (-O3 -fno-tree-vectorize): 0 loops vectorized
\end{lstlisting}
\qty{32}{\byte} vectors are \num{8} float32 lanes (native), \qty{16}{\byte} are \num{4} (autovec);
GCC's ``unroll factor'' is the vectorization factor, not extra unrolling. The native epilogue is one
\qty{128}{bit} step (if \num{4} or more elements remain) plus up to three scalar iterations;
autovec's is one \qty{8}{\byte} step plus at most one scalar iteration. The loop is
\emph{versioned} on a run-time aliasing check between \code{x} and \code{y}, and GCC does not peel
iterations to reach an aligned start.

\subsection{Disassembly}
The native SAXPY inner loop is one fused multiply--add on \code{ymm} registers:
\begin{lstlisting}[style=excerpt]
vmovups   (%rdi,%rax,1),%ymm0           # load 8 x[i]
vfmadd213ps (%rsi,%rax,1),%ymm1,%ymm0   # ymm0 = a*x + y   (one FMA, 8 lanes)
vmovups   %ymm0,(%rsi,%rax,1)           # store 8 y[i]
\end{lstlisting}
The scalar build handles one element at a time and cannot fuse, because the generic \code{x86-64}
target lacks FMA:
\begin{lstlisting}[style=excerpt]
movss (%rdi,%rax,4),%xmm0  # x[i]
mulss %xmm1,%xmm0          # a*x[i]
addss (%rsi,%rax,4),%xmm0  # + y[i]   (separate mul + add: generic x86-64 has no FMA)
movss %xmm0,(%rsi,%rax,4)  # y[i]
\end{lstlisting}

The dot product is the exception: native multiplies eight lanes at once but sums them one lane at
a time, in source order:
\begin{lstlisting}[style=excerpt]
vmulps (%rdi,%rax,1),%ymm1,%ymm1  # 8 products at once ...
vaddss %xmm1,%xmm0,%xmm0          # ... but summed one lane at a time,
vaddss %xmm3,%xmm0,%xmm0          #     in strict source order
...
\end{lstlisting}

\subsection{Stride and gather code}
\code{saxpy_f32_stride} shows \code{ymm} instructions in \code{native} only because GCC versions
the loop on \code{s == 1}; for $s>1$ it runs a scalar \code{vfmadd213ss} loop (\code{mulss}/\code{addss}
in the other builds). No build vectorizes the plain gather (\enquote{no vectype for stmt}), so we
added an explicit AVX2 version, \code{saxpy_f32_gather_avx2}, compiled only into \code{native}:
\begin{lstlisting}[style=excerpt]
vmovdqu    (%rax),%ymm1                 # 8 indices idx[i..i+7]
vgatherdps %ymm5,(%rdi,%ymm1,4),%ymm0   # gather 8 x[idx]
vgatherdps %ymm7,(%rsi,%ymm1,4),%ymm3   # gather 8 y[idx]
vfmadd132ps %ymm4,%ymm3,%ymm0           # one 8-lane FMA
vmovss     %xmm0,(%rsi,%rdx,4)          # ... then 8 scalar stores:
vextractps $0x1,%xmm0,(%rsi,%rdx,4)     #     AVX2 has no scatter
\end{lstlisting}
The loads are SIMD; the stores cannot be, since AVX2 has no scatter.

\subsection{Interpretation}
The evidence matches the hypothesis: \code{native} uses \code{ymm} for SAXPY, multiply and the
stencil (with FMA in SAXPY and the stencil), \code{autovec} uses \code{xmm}, and \code{scalar} is
genuinely scalar. Tails take a narrower vector step plus a scalar remainder; GCC does not use AVX2's
masked \code{vmaskmovps} (no AVX-512 masks here). The dot product's packed multiplies feed a
latency-bound scalar accumulator, so ``vectorized'' does not mean ``fast'' (\cref{sec:negative}).
